\section{Introduction}

An odd cycle in the permutation induced by two lines of a Latin square
supports a simple trade: swap the two lines on that cycle.  The move preserves
the Latin property.  In the row and column views it reverses the conventional
row--column Latin sign; in the symbol view it preserves that sign and reverses
the two parastrophically transported signs.  This local operation is therefore
a natural primitive in parity-based approaches related to the Alon--Tarsi
setting, but its three sign conventions must be kept distinct.  B\'erczi's
experimental program studies such local corrections as ingredients for
global constructions~\cite{Berczi2026}; it does not supply the theorems proved
here.

We isolate the underlying existence question.  For each of the row, column,
and symbol views, ask whether some pair of lines induces a permutation with a
nontrivial odd cycle.  The three answers form a pattern in $\{F,T\}^3$.  A
square is FFF when all three answers are negative.  Equivalently, every union
of two link factors in the tripartite representation has component lengths
divisible by four.

Our first contribution is structural.  Every odd-order Latin square is TTT,
whereas the Cayley table of a finite group is FFF exactly when the group is a
$2$-group.  A basis-family closure criterion detects group isotopy.  We then
show that failure patterns are monotone from Latin subsquares and analyze
quasigroup congruences: odd block size forces TTT, while two-element blocks
preserve the quotient pattern exactly.  Thus a hypothetical FFF quasigroup of
order $2p$, for odd prime $p$, must be congruence-simple in every isotope.
This excludes all congruence-imprimitive constructions at order $10$ but does
not by itself decide the simple case.

The remaining binary line-incidence defect also has an exact subsquare
meaning.  Nonuniversal vectors in the left binary line kernel are in
bijection with half-order subsquares, and the quotient classes group them
into complementary four-packs.  This yields a closed counting formula and
forces maximal binary line rank for every FFF order congruent to $2$ modulo
$4$ and greater than $2$.  At order $10$ this is necessary, not sufficient.

More generally, Moorhouse's $3$-net character formula specializes as follows:
a nonuniversal line dependency over a prime field is a
balanced cyclic quotient homotopy, and after isotopy its fibres form a
quasigroup congruence.  Consequently every Smith torsion prime divides the
order.  For order $2p$, with $p$ odd prime, the congruence obstruction shows
that even one F view forces a saturated integer line-incidence lattice.

The group construction is not the whole phenomenon.  A twisted binary
extension of $C_4$ gives an explicit non-group-isotopic FFF loop of order $8$.
Moreover, direct products satisfy the exact componentwise formula
\[
  \pat(L\times K)=\pat(L)\mathbin{\mathrm{OR}}\pat(K),
\]
and a product is group-isotopic exactly when both factors are.  These theorems
turn one order-$8$ example into infinite power-of-two families.

Our computational contribution is exhaustive.  McKay, Meynert, and Myrvold
enumerated the main classes of small Latin squares~\cite{McKayMeynertMyrvold2007};
the associated ANU file contains one representative of each of the $283657$
order-$8$ main classes~\cite{McKayLatinData}.  Our scan finds exactly $230$
FFF main classes.  Only $5$ are group-isotopic and $225$ are not, and all
eight ordered failure patterns occur.  Independently, fresh exact-cover
generation of all $9408$ reduced order-$6$ squares finds no FFF square and
matches the official ANU set exactly.

Combining theory and computation proves that every order $2^m$, $m\geq3$,
realizes all eight patterns and contains a non-group-isotopic FFF square.  We
do not claim that the product construction is injective on isotopy or main
classes.  We also do not prove the Alon--Tarsi conjecture.  The accompanying
complete research report proves by exhaustive master searches that order $10$
has no FFF square; that computation is outside this compact paper's theorem
scope.  The resulting open conjecture is narrower: nontrivial FFF Latin
squares may exist exactly in power-of-two orders.

The longer companion research report in the repository develops additional
cycle-parity, subset-action, primitive-group, and flag-surface machinery
and proves nonexistence at order $10$.  None of that machinery is needed for
the results in this compact paper.
